\section*{Chapter \ref{ch:g11:gates-to-arithmetic-unit} --- From Gates to an Arithmetic Unit}

\begin{solution}{exo:g11:gates-to-arithmetic-unit:1}
$p = 1 \oplus 0 = 1$ and $g = 1 \wedge 0 = 0$; then $s = p \oplus c = 0$ and the second carry is $p \wedge c = 1$, so
$c' = 0 \vee 1 = 1$. Indeed $1 + 0 + 1 = 2$, written $10$.
\end{solution}

\begin{solution}{exo:g11:gates-to-arithmetic-unit:2}
$c_0 = 0$. Column $0$: $1 + 0 + 0 = 1$, so $s_0 = 1$, $c_1 = 0$. Column $1$: $1 + 1 + 0 = 2$, so $s_1 = 0$,
$c_2 = 1$. Column $2$: $1 + 1 + 1 = 3$, so $s_2 = 1$, $c_3 = 1$. Column $3$: $0 + 0 + 1 = 1$, so $s_3 = 1$,
$c_4 = 0$. The sum is $1101$, that is $7 + 6 = 13$.
\end{solution}

\begin{solution}{exo:g11:gates-to-arithmetic-unit:3}
$s_1 s_0 = 00$ gives $x_0 = 1$, $01$ gives $x_1 = 0$, $10$ gives $x_2 = 0$, $11$ gives $x_3 = 1$.
\end{solution}

\begin{solution}{exo:g11:gates-to-arithmetic-unit:4}
$2^3 = 8$ outputs, numbered $0$ to $7$; for $101$ the output number $5$ is $1$ and the seven others are $0$.
\end{solution}

\begin{solution}{exo:g11:gates-to-arithmetic-unit:5}
The carry is \code{AND(a, b)}; the sum is $1$ when at least one input is $1$ but not both, so it is
\code{AND(OR(a, b), NOT(AND(a, b)))}. The file \code{exercises/circuits.py} has it, and the tests check the four
rows.
\end{solution}

\begin{solution}{exo:g11:gates-to-arithmetic-unit:6}
Six rows are wrong: every row except $000$ and $001$. With an AND, the first ``sum'' is $1$ only when $a = b = 1$,
so it is wrong whenever exactly one of $a$, $b$ is $1$ (four rows); and when $a = b = 1$, the two carries are both
$1$ and the sum bit becomes $c$ instead of $1 \oplus 1 \oplus c$, wrong in both rows. The \code{assert} of
\code{full_adder.py}, which tries all eight rows, stops at the third, $010$.
\end{solution}

\begin{solution}{exo:g11:gates-to-arithmetic-unit:7}
Split \code{sel} into its two bits, then choose within each pair and between the pairs:

\smallskip\noindent\hspace*{2em}\code{s1, s0 = (sel >> 1) & 1, sel & 1}

\noindent\hspace*{2em}\code{low = mux2(s0, inputs[0], inputs[1])}

\noindent\hspace*{2em}\code{high = mux2(s0, inputs[2], inputs[3])}

\smallskip\noindent and the result is \code{mux2(s1, low, high)}. The file \code{exercises/circuits.py} writes the
three calls as one expression.
\end{solution}

\begin{solution}{exo:g11:gates-to-arithmetic-unit:8}
Five gates per \omterm{def:g11:gates-to-arithmetic-unit:full-adder}{full adder}: $5n$. With a half adder (two gates) in the first column: $5(n - 1) + 2 = 5n - 3$.
\end{solution}

\begin{solution}{exo:g11:gates-to-arithmetic-unit:9}
From $c_i$ to $c_{i+1}$ the signal crosses the AND of the second half adder and the OR: $2$ gates. In column $0$, the
carry $c_1$ needs the XOR, then an AND, then the OR: level $3$. Each later column adds $2$, so $c_n$ is ready at level
$3 + 2(n - 1) = 2n + 1$, and that is the longest chain. For $64$ bits, $129$ gates.
\end{solution}

\begin{solution}{exo:g11:gates-to-arithmetic-unit:10}
It computes $a + \overline{b} = a + (2^n - 1 - b)$, that is $a - b - 1$ modulo $2^n$: $100 - 27$ gives $72$
instead of $73$. Every result is one too small.
\end{solution}

\begin{solution}{exo:g11:gates-to-arithmetic-unit:11}
For ADD, the result is \texttt{(x + y) \% 256} and $C$ is \code{int(x + y > 255)}; for SUB, \texttt{(x - y) \% 256} and
\code{int(x < y)}; $Z$ is \code{int(r == 0)} and $N$ is \code{r >> 7} (\code{exercises/circuits.py}). Two versions
written in very different ways rarely make the same mistake: comparing them on many inputs finds the mistakes of
either. The tests compare the gate ALU with this one and with the toy machine.
\end{solution}

\begin{solution}{exo:g11:gates-to-arithmetic-unit:12}
The table shows $s = 1$ exactly on the rows with one or three ones, an odd number; $a \oplus b$ is $1$ when $a + b$ is
odd, and XOR with $c$ flips that when $c = 1$, so $a \oplus b \oplus c$ is $1$ exactly when $a + b + c$ is odd:
$s = a \oplus b \oplus c$. The table shows $c' = 1$ exactly on the rows $011$, $101$, $110$, $111$, those with at
least two ones. The term $a \wedge b$ is $1$ on $110$ and $111$; the term $c \wedge (a \oplus b)$ on $011$ and $101$;
their OR is $1$ on these four rows and no other. The three terms $a \wedge b$, $b \wedge c$, $a \wedge c$ say ``these
two are $1$'', and their OR is $1$ exactly when at least two inputs are: the same four rows.
\end{solution}

\begin{solution}{exo:g11:gates-to-arithmetic-unit:13}
Bit $i$ of $\overline{b}$ is $1 - b_i$, so $\overline{b} = \sum_i (1 - b_i) 2^i = (2^n - 1) - b$. Then
$a + \overline{b} + 1 = a - b + 2^n$, which is $a - b$ modulo $2^n$. The adder keeps $n$ bits, and its carry out is
$1$ exactly when the true sum $a - b + 2^n$ is at least $2^n$, that is when $a \ge b$.
\end{solution}

\begin{solution}{exo:g11:gates-to-arithmetic-unit:14}
Subtract $y$ from $x$. If the result is $0$, \code{"="}; otherwise, a carry out of $0$ means a borrow, so $x < y$ and
the answer is \code{"<"}; a carry out of $1$ means $x > y$ (\code{exercises/circuits.py}). These are the \omterm{def:g11:gates-to-arithmetic-unit:flag}{flags} $Z$
and $C$ of the instruction \code{CMP} of the toy machine.
\end{solution}

\begin{solution}{exo:g11:gates-to-arithmetic-unit:15}
With $S = R = 1$, from $(Q, \overline{Q}) = (0, 1)$: both NORs see a $1$, so the next state is $(0, 0)$, and it stays.
Releasing both at once, each gate sees two zeros: $(1, 1)$; then each sees a $1$: $(0, 0)$; and so on for ever, as
the output of \code{latch.py} shows. If the gate of $Q$ is faster, it switches first to $Q = \text{NOR}(0, 0) = 1$;
the other gate then sees $Q = 1$ and stays at $0$: the \omterm{def:g11:gates-to-arithmetic-unit:latch}{latch} ends at $Q = 1$. If the other gate is faster, it ends at
$Q = 0$. The final state depends on which gate is faster, which nobody controls; and while $S = R = 1$, $\overline{Q}$
is not the opposite of $Q$. Both are reasons to forbid these inputs.
\end{solution}

\begin{solution}{pb:g11:gates-to-arithmetic-unit:1}
\textbf{1.} The level of its output: one more than the latest of its inputs, since a gate can only switch once all
its inputs are there.

\textbf{2.} \code{p} is the XOR of $a$ and $b$, \code{s} the XOR of \code{p} and $c$, then two ANDs, of $a$ and $b$
and of \code{p} and $c$, and the OR of these: five gates, as in the figure.

\textbf{3.} $p$ at $1$, $s$ at $2$; the ANDs at $1$ and $2$, the carry out at $3$.

\textbf{4.} $p$ at $1$, $s$ at $k + 1$; the AND of $a$ and $b$ at $1$, the AND of $p$ and $c$ at $k + 1$, and the
carry out at $k + 2$.

\textbf{5.} By question 3, $c_1$ is at level $3$; by question 4, a carry at level $k$ gives the next one at $k + 2$.
So $c_1, c_2, c_3, \dots$ are at $3, 5, 7, \dots$: $c_i$ is at $2i + 1$.

\textbf{6.} The last sum bit $s_{n-1}$ is at $(2(n - 1) + 1) + 1 = 2n$, the carry out $c_n$ at $2n + 1$. The output
gives $5$ for $2$ bits and $129$ for $64$: $2n + 1$, the carry out being the last signal ready.

\textbf{7.} $5n$: $40$, $80$, $160$ and $320$ gates.

\textbf{8.} $129 \times 10$~ps $= 1.29$~ns, so about $775$ million additions in a row per second.

\textbf{9.} The carry chain: each sum bit is ready one level after its carry, so the sums wait for the carries.

\textbf{10.} Its inputs, the high halves of $A$ and $B$, are there from the start; only its carry in is missing, and
that carry can only be $0$ or $1$.

\textbf{11.} Once assuming a carry of $0$, once assuming $1$; when the low half's carry arrives, a \omterm{def:g11:gates-to-arithmetic-unit:mux}{multiplexer} per
output bit keeps the copy made with the right assumption.

\textbf{12.} One \omterm{def:g11:gates-to-arithmetic-unit:mux}{multiplexer} per sum bit of the high half and one for the carry out: $n/2 + 1$. Each takes two ANDs
and an OR, and the NOT of the select signal is shared by all.

\textbf{13.} Three ripple adders of $n/2$ bits, $3 \times 5n/2$ gates, plus $3(n/2 + 1)$ for the \omterm{def:g11:gates-to-arithmetic-unit:mux}{multiplexers} and $1$
NOT: $15n/2 + 3n/2 + 4 = 9n + 4$. For $n = 64$, $580$, as printed.

\textbf{14.} The low half is a ripple adder of $n/2$ bits: its carry is at level $2(n/2) + 1 = n + 1$. The high halves
are ready earlier, at level $n + 1$ at most. The NOT of the carry is at $n + 2$, the AND that uses it at $n + 3$, the
OR at $n + 4$.

\textbf{15.} $64 + 4 = 68$.

\textbf{16.} $580/320 \approx 1.81$ times the gates, for $68/129 \approx 0.53$ times the levels.

\textbf{17.} The first block of $16$ bits gives its carry at level $33$. The three other blocks are computed both ways
in parallel, and ready at level $33$ too. Each \omterm{def:g11:gates-to-arithmetic-unit:mux}{multiplexer} stage adds $3$ levels to the chosen carry: $36$ after the
second block, $39$ after the third, and the outputs of the last block at about $42$. The gates grow again, to about
$5 \cdot 16 + 3 \cdot 2 \cdot 5 \cdot 16 + 3 \cdot 3 \cdot 17 + 3 \approx 716$.

\textbf{18.} Every instruction must finish within one beat of the clock, so the slowest circuit sets the speed of the
whole processor. Gates are tiny and cheap, counted in billions on one chip; time is not.

\textbf{19.} More blocks, and blocks inside blocks, reduce the depth further: adders used in processors make the
depth grow much more slowly than the number of bits, so that doubling the bits adds only a few levels. The shortest
imaginable adder writes each output as one OR of ANDs of the inputs, two levels deep, but that needs gates with dozens
of inputs, which do not exist; built from two-input gates they bring the levels back. Speed always costs gates.

\textbf{20.} A $64$-bit \omterm{def:g11:gates-to-arithmetic-unit:ripple}{ripple-carry adder} needs $320$ gates, five per \omterm{def:g11:gates-to-arithmetic-unit:full-adder}{full adder}, and its longest chain crosses
$129$ gates, the carry gaining two levels per column. A carry-select adder of two halves needs $580$ gates, $1.81$
times as many, and $68$ levels, $0.53$ times as many: it nearly halves the delay for less than twice the gates.
\end{solution}
